% [original page: 72]
מוכח, שמקיום המשוואה
\[ ax^2 - by^2 = \pm 2^2 \, , \quad (x,y) = 1 \, , \]
נובעת קיומה של המשוואה
\[ as^2 - bt^2 = \pm 1 \]
והיפוכו של דבר זה נכון בתנאי, שיש פתרון למשוואה של
\[ u^2 - abv^2 = 4 . \]
בתנאי: $(u,v) = 1$.

\medskip
\noindent הכללה נוספת של משפט ריכלדה - פאלמשטרים גוררת אחריה את הנוסחאות
\[ \begin{cases} ax_n^2 + by_n^2 = 2u_{2n-1} \\ x_n y_n = v_{2n-1} \end{cases} \]
בשביל כל פתרונות המשוואה
\[ ax^2 - by^2 = \pm 4 \]
כשיש לה בכלל פתרון, במקום ש
\[ u^2 - abv^2 = 4 \]
ו
\[ \frac{u_k + v_k\sqrt{ab}}{2} = \pm \left(\frac{u_1 \pm v_1\sqrt{ab}}{2}\right)^k \, , \quad k>0 \]
מכאן מתקבלות מיד הנוסחאות:
\[ \frac{x_n\sqrt{a} + y_n\sqrt{b}}{2} = \pm \left(\frac{x_1\sqrt{a} \pm y_1\sqrt{b}}{2}\right)^{2n - 1} \]
כאשר $x_1$, $y_1$ הנו הפתרון המינימלי של המשוואה הנ"ל בשביל כל
\[ n = 1 , 2 , 3 , \ldots \]
כן קיים המשפט, שאם יש פתרון למשוואה
\[ u^2 - abv^2 = 4 \]
\[ (u , v ) = 1 \]
אזי יש פתרון לאחת המשוואות מהצורה

% [original page: 73]
\noindent בתנאי
\[ a > 1 \]
רק לאחת מבין אלה.

\medskip
\noindent כהכללת הקשר של פואנקרה אפשר לראות את הקשר:
\[ s_k\sqrt{a} + t_k\sqrt{b} = \left(\frac{x_k\sqrt{a} + y_k\sqrt{b}}{2}\right)^3 = \left(\frac{x_3\sqrt{a} + y_3\sqrt{b}}{2}\right)^k \]
במקום ש $x$, $y$ ; $s$, $t$ הם פתרונות המשוואות
\[ \begin{cases} ax^2 - by^2 = \pm 4 \\ as^2 - bt^2 = \pm 1 \end{cases} \]
בהתאמה.

\medskip
\noindent בכל המסקנות הללו יש לראות את תכנית הכללת משוואת פל - פרמה, בכיוון אחד; כלומר את הכללת המשוואה ביחס למקדמיה. אולם הכללתה נעשתה גם בכיוון אחר, והוא ביחס למעריכים. אבל אל נושא זה נשוב בסעיף המיוחד לכך.

\bigskip
\noindent\textbf{2.} אנו עוברים עתה לטיפול במשוואה
\[ ax^2 - by^2 = \pm 1 \tag{1} \]
בתנאים: $a>1$, $b>1$

\noindent מטרתנו לקבוע שיטה מושלמת שעל פיה אפשר לדעת אם יש פתרון למשוואה ובמקרה שיש - למצוא את פתרונותיה.

\noindent על ידי העלאת המשוואה (1) לריבוע אנו מקבלים משוואה מצורת משוואת פל
\[ \mathcal{E}^2 - ab\eta^2 = 1 \tag{2} \]
וקיים
\[ ax^2 + by^2 = \mathcal{E} \, , \quad 2xy = \eta \tag{3} \]
בשביל המספרים $\mathcal{E}$, $\eta$ השלמים, כפי שרואים מיד.

% [original page: 74]
\noindent מכאן נובע, שבשביל כל פתרון $x$, $y$ של המשוואה (1) מתאים פתרון $\mathcal{E}, \eta$ מסויים של משוואה (2) באמצעות הנוסחאות (3).

\noindent אולם התאמה זאת אינה ניתנת גם להיפוך; כי כפי שאפשר להיווכח אין למצוא פתרון $x$, $y$ למשוואה (1) בשביל כל פתרון $\mathcal{E}, \eta$ של המשוואה (2) מתוך הנוסחאות (3) במספרים $x$, $y$ שלמים.

\noindent עלינו לבחור, איפוא, את אותם הפתרונות $\mathcal{E}, \eta$ של משוואה (2) אשר בשבילם כן אפשר להתאים מספרים $x$, $y$ שלמים המקיימים את המשוואה (1).

\medskip
\noindent כשהמספרים $x$, $y$ הם שלמים נקבל מיד
\[ (ax^2 + by^2)^2 - ab(2xy)^2 = (ax^2 - by^2)^2 = 1 \]
ולכן
\[ ax^2 - by^2 = \pm 1 \]
כלומר: כל המספרים $x$, $y$ השלמים המקיימים את הנוסחאות (3) קובעים אחד הפתרונות של המשוואה (1).

\medskip
\noindent נצטמצם, איפוא, לבעיה זאת האחרונה. מהנוסחאות (3)(1) אנו מקבלים את הקונגרואנציות
\[ \mathcal{E} \pm 1 \equiv 0 \pmod{2a} \tag{4} \]
\[ \mathcal{E} \mp 1 \equiv 0 \pmod{2b} \]
ולכן נובע, שהתנאי ההכרחי לקיום המשוואות (2)(3) הוא קיום היחסים האחרונים המתבטאים ע"י שתי הקונגרואנציות.

\noindent אולם בנקל ניווכח, שקיום שתי הקונגרואנציות הוא גם תנאי מספיק למציאות פתרונות עפ"י הדרישות (2)(3).

\noindent ואמנם את המשוואה (2) אפשר לכתוב בצורה:
\[ (\mathcal{E} + 1)(\mathcal{E} - 1) = ab\eta^2 . \]
נשים עתה לב לעובדה:
\[ (\mathcal{E} + 1, \mathcal{E} - 1) = (\mathcal{E} + 1, 2) \leq 2 \]

% [original page: 75]
\noindent מהפירוק החד-ערכי הקיים לגבי מספרים רציונליים שלמים נסיק מכאן מיד, כי:
\[ \mathcal{E} \pm 1 = 2ap^2 \, , \quad \mathcal{E} \mp 1 = 2bq^2 \]
בשביל מספרים $p$, $q$ שלמים מסויימים.

\noindent אולם ערכים אלו מקיימים, כפי שרואים מיד, את הקשרים:
\[ 2pq = \eta \, , \quad ap^2 + bq^2 = \mathcal{E} \]
מטעמי סימטריות הפתרונות $x$, $y$ של המשוואה (1) ביחס לציר ה-$X$ ולציר ה-$Y$ נצטמצם להבא, מבלי להגביל את הכלליות, רק בחקירת הערכים $\eta$, $\mathcal{E}$ הטבעיים, וביחד עם זאת נצטמצם, איפוא, לערכים $\mathcal{E}$, $\eta$ טבעיים.
% ILLEGIBLE: this sentence is partly obscured by dot-matrix noise; approximate reading given, exact wording of the middle clause uncertain.

\medskip
\noindent מהנוסחה הכללית לפתרונות משוואת פל ניתן מיד לאפשר לכתוב
\[ \mathcal{E}_n = \mathcal{E}_1^n + abM_n \tag{5} \]
במקום ש $M_n$ הנו מספר שלם מסויים המקיים את הנוסחה:
\[ M_n = \frac{n(n-1)}{1 \cdot 2}\mathcal{E}_1^{n-2}\eta_1^2 + \eta_1 \, abK \tag{6} \]
לגבי $K$ שלם מסויים.

\noindent נניח עתה כי:

\medskip
\noindent\underline{$n$ = למספר אי-זוגי.}

\noindent אם מתקיימות הנוסחאות (2)(3) בשביל $\mathcal{E}_n, \eta_n$ אזי לפי משפט העזר 1, נסיק מיד ובייחד עם זה כל הערכים $\eta$ מוכרחים להיות זוגיים.

\noindent לפיכך נובע על סמך (4)(3), כי:
\[ \begin{cases} \mathcal{E}_n + 1 = P(\mathcal{E}_1 + 1) + 2abr \\ \mathcal{E}_n - 1 = Q(\mathcal{E}_1 - 1) + 2abr \end{cases} \]
בשביל מספרים $P$, $Q$, $r$ שלמים; על סמך זה נוכל לקבוע כי:
\[ \left. \begin{aligned} \mathcal{E} \pm 1 &\equiv 0 \pmod{2a} \\ \mathcal{E} \mp 1 &\equiv 0 \pmod{2b} \end{aligned} \right\} \tag{7} \]

% [original page: 76]
\[ \mathcal{E}_{2n-1} \pm 1 \equiv 0 \pmod{2a} \]
\[ \mathcal{E}_{2n-1} \mp 1 \equiv 0 \pmod{2b} \]
בשביל כל
\[ n = 1 , 2 , 3 , \ldots \]
אני טוען, שגם ההיפוך נכון; כלומר:
\[ n > 1 , \; \begin{cases} \mathcal{E}_{2n-1} \pm 1 \equiv 0 \pmod{2a} \\ \mathcal{E}_{2n-1} \mp 1 \equiv 0 \pmod{2b} \end{cases} \]
\[ \begin{cases} \mathcal{E}_1 \pm 1 \equiv 0 \pmod{2a} \\ \mathcal{E}_1 \mp 1 \equiv 0 \pmod{2b} \end{cases} \tag{8} \]
כי אלמלא כן, היינו מסיקים על סמך הקשר
\[ (\mathcal{E}_1 + 1)(\mathcal{E}_1 - 1) = ab\eta_1^2 \]
שקיים
\[ \mathcal{E}_1 \pm 1 \equiv 0 \pmod{2d} \]
\[ \mathcal{E}_1 \mp 1 \equiv 0 \pmod{2d_1} \]
במקום ש
\[ dd_1 = ab \, , \quad d \neq a \, , \quad d_1 \neq b \, ; \]
כי לפי משפט העזר 1 ולפי ש-$\eta_1$ נובע ש הנן זוגי.

\noindent אבל לפי"ז היינו מסיקים, לפי (7), שמתקיימות ג"כ:
\[ \mathcal{E}_{2n-1} \pm 1 \equiv 0 \pmod{2d} \]
\[ \mathcal{E}_{2n-1} \mp 1 \equiv 0 \pmod{2d_1} \]
ודבר זה לא יתכן, ביחד עם הקונגרואנציות הקיימות לפי הנחתנו, מכיוון שאין כל מחלק משותף $>2$ לשני המספרים $\mathcal{E}-1$, $\mathcal{E}+1$.

% [original page: 77]
\noindent יהי עתה

\medskip
\noindent\underline{$n$ = למספר זוגי.}

\noindent במקרה זה מתחלק $\mathcal{E}_1^n - 1$ ב-$\mathcal{E}_1^2 - 1$ ולכן אפשר לכתוב, בהתאם לנוסחה (5):
\[ \mathcal{E}_n - 1 = N(\mathcal{E}_1^2 - 1) + abM_n \tag{9} \]
אבל $\mathcal{E}_1^2 - 1$ מתחלק ב-$ab$, כפי שרואים מיד מנוסחה (2); חוץ מזה מתקיים:
\[ (\mathcal{E}_n + 1 \, , \, \mathcal{E}_n - 1) \leq 2 \]
ולכן נוכל להחליט מיד, שהקונגרואנציות (4) יכולות להתקיים לכל היותר, בצורה:
\[ \begin{cases} \mathcal{E}_n + 1 \equiv 0 \pmod{2^2} \\ \mathcal{E}_n - 1 \equiv 0 \pmod{2B} \end{cases} \tag{10} \]
במקום:
\[ 2B = ab \, . \]
כלומר: כאשר המשוואה (1) מקבלת את הצורה:
\[ 2x^2 - By^2 = 1 \, . \tag{11} \]
מכאן:

\medskip
\noindent\underline{משפט 8.}

\noindent א) בשביל כל פתרונות המשוואה
\[ \mathcal{E}^2 - D\eta^2 = 1 \, , \quad D = ab \]
קיימים:
\[ 2(\mathcal{E}_{2n} + 1) = u^2 \]
\[ 2(\mathcal{E}_{2n} - 1) = Dv^2 \]
% NOTE: a marginal/interlinear notation "u | \eta_{2n}" appears here in the source (small superscript-like mark to the right of the displayed equations); interpreted per the divisibility convention as $u \mid \eta_{2n}$, but the exact placement is uncertain — see transcription notes.
\[ u \mid \eta_{2n} \]
בשביל כל $n$ טבעי.

\noindent ב) אם קיימות הנוסחאות:

% [original page: 78]
\[ 2(\mathcal{E}_{2n-1} - 1) = bq_n^2 \]
\[ 2(\mathcal{E}_{2n-1} + 1) = ap_n^2 \qquad , \; a \neq 2 \]
\[ p_n \mid 2\eta_{2n-1} \, , \qquad q_n \mid 2\eta_{2n-1} \]
% NOTE: the two divisibility relations above are transcribed from a fraction-like stacked notation in the source ("p_n" over "2\eta_{2n-1}"), per the style guide's convention that such stacked layouts denote divisibility in this number-theoretic context.
בשביל $n$ מסויים, אזי הן נשארות בתקפן בשביל כל
\[ n = 1 , 2 , 3 , \ldots \]
כאשר
\[ ab = D. \]
במשפט זה אפשר לראות הכללה של משפט ריכלדה ו-פאלמשטרים (ראה: \LR{L.E. Dickson, \textit{History of the Theory of Numbers}, vol. II, pp. 392}).

\noindent המתימטיקאים הללו קבעו את המשפט הבא:

\begin{quote}
"הביטוי $\sqrt{P_n}$, כאשר $P_n = 2(\mathcal{E}_{2n} + 1)$, הנו ריבוע שלם, והמספר $\eta_{2n}$ מתחלק ב-$\sqrt{P_n}$; אם הביטוי $Q_n = 2(\mathcal{E}_{2n-1} - 1)$ הנו ריבוע שלם בשביל $n$ מסויים, אזי הוא גם ריבוע שלם בשביל כל $n$, והמספר $\eta_{2n-1}$ מתחלק ב-$\sqrt{Q_n}$."
\end{quote}

\noindent השוואה קלה של המשפטים האחרונים מעמידה אותנו על צדקת טענתנו.

\noindent כמסקנה ישירה ממשפט 8 אנו יכולים לקבוע את

\medskip
\noindent\underline{משפט 9.}

\noindent כל הערכים המקיימים את הסימונים (1)(2)(3), אם ישנם כאלה בכלל, נמצאים בנוסחאות:
\[ \begin{cases} ax_n^2 + by_n^2 = \mathcal{E}_{2n-1} \\ 2x_ny_n = \eta_{2n-1} \end{cases} \qquad n = 1 , 2 , 3 , \ldots \tag{12} \]

% [original page: 79]
\noindent במשפט זה גמרנו את חקירת המשוואה (1) בתנאי:
\[ a \neq 2 . \]
נשאר לנו עוד לסיים את חקירת המקרה המיוחד, שבו $a=2$, כלומר: לברר את כל מה שנוגע למשוואה מהצורה (11).

\medskip
\noindent מהקונגרואנציות (10) נוכל להסיק על סמך העובדה
\[ (\mathcal{E}_n - 1 \, , \, \mathcal{E}_n + 1) \leq 2 \, , \]
ש-$\mathcal{E}_n - 1$ אינו יכול להתחלק ב-4.

\noindent אולם כאשר $ab$ הנו זוגי מוכרח להיות $\mathcal{E}_1$ אי-זוגי, כפי שרואים מיד ממשוואה (1) עצמה, לפיכך מתחלק הביטוי
\[ \mathcal{E}_1^2 - 1 = (\mathcal{E}_1 + 1)(\mathcal{E}_1 - 1) \]
ב-4.

\noindent מ-(9) נוכל, איפוא, להסיק שהתנאי
\[ \mathcal{E}_n - 1 \not\equiv 0 \pmod{4} \]
יכול להתקיים אז ורק אז כאשר $M_n$ הנו אי-זוגי. כי $B$ מוכרח להיות אי-זוגי, כפי שנובע מהמשוואה (11) עצמה, ולכן
\[ ab = 2B \equiv 2 \pmod{4} \]
אבל מנוסחה (6) אנו מסיקים ש-$M_n$ יכול להיות אי-זוגי אז ורק אז כאשר $\eta_1$ הנו אי-זוגי ו-$n$ אינו מתחלק ב-4, מכיוון ש-$ab$ הנו זוגי, ו-$\mathcal{E}_1$ הנו אי-זוגי.

\noindent מאידך, נוכל להיווכח שהתנאים הללו ניתנים להיפוך; כלומר: אם מתקיימים התנאים האחרונים, אזי מתקיימות הקונגרואנציות (10).

\noindent הלאה, מהשוויון
\[ (\mathcal{E}_n + 1)(\mathcal{E}_n - 1) = ab\eta_n^2 \]
אנו מחליטים מיד בתנאים אלו, שאחד המספרים
\[ \mathcal{E}_n + 1 \, , \quad \mathcal{E}_n - 1 \]
מתחלק ב-4, ובעונה אחת השני הנו זוגי, מכיוון ש-$ab\eta_n^2$ מתחלק ב-8.

\noindent אולם בתנאים שלנו אין הביטוי $\mathcal{E}_n - 1$ מתחלק ב-4, כפי שכבר בררנו,

% [original page: 80]
\noindent ומאידך הוא מתחלק ב-2 ($ab = 2B$) ולכן מתחלק $\mathcal{E}_n + 1$ ב-4.

\noindent מהשיקולים הללו אנו מגיעים לידי

\medskip
\noindent\underline{משפט 10.}

\noindent א) כאשר
\[ \eta_1 \equiv 1 \pmod{2} \]
יש פתרון למשוואה
\[ 2x^2 - By^2 = 1 \]
וכל הערכים $(x,y)$ המקיימים את המשוואה האחרונה נמצאים בנוסחאות:
\[ \begin{cases} 2x_n^2 + By_n^2 = \mathcal{E}_{2(2n-1)} \\ 2x_ny_n = \eta_{2(2n-1)} \end{cases} \qquad n = 1 , 2 , 3 , \ldots \tag{12} \]
ב) כאשר
\[ \eta_1 \not\equiv 1 \pmod{2} \]
יש לנהוג בפתרון המשוואה עפ"י השיטה הנכללת במשפט 9.

\medskip
\noindent בזה נגמרה חקירת המשוואה
\[ ax^2 - by^2 = 1 \]
בכל המקרים.

\bigskip
\noindent\textbf{3.} נביא עתה כמה מהמשפטים בקשר למשוואה שחקרנו, אשר מהם נעשה עוד שימוש בתורת התבניות הדו-צדדיות.

\noindent קל להיווכח, שאין מקום לקונגרואנציות (4) בשביל שני זוגות מודולים $(a,b)$, $(a_1,b_1)$ בבת-אחת אם נבחר במקום $\mathcal{E}$ את $\mathcal{E}_1$.

\noindent מכאן נובע מיד, על סמך כל השיקולים שבסעיף הקודם,

\medskip
\noindent\underline{משפט 11.}

\noindent למשוואה מהצורה (1) יש פתרון לכל היותר בשביל זוג גורמים $a,b$ אחד של המספר $D$, זרים ביניהם, בתנאי: $a>1$, $b>1$.

% [original page: 81]
\noindent כמסקנה מכאן אנו מקבלים במישרים את

\medskip
\noindent\underline{משפט 12.} משוואה (1) אינה יכולה להתקיים בשביל שני היימנים בבת-אחת.

\medskip
\noindent נקח עתה לתשומת לבנו, שהמשוואה
\[ x^2 - Dy^2 = -1 \]
יכולה להיות מקרה פרטי של המשוואה (1), ונביא בחשבון את המסקנות מתורת שברי-הטרסרת לגבי משוואת פל ומיד יתקבל

\medskip
\noindent\underline{משפט 13.} אם המחזור שבפיתוח $\sqrt{ab}$ לשבר משולב מרכב ממספר אי-זוגי של אברים אזי, אין פתרון למשוואה (1) בשביל
\[ a>1 , \quad b>1 . \]
לפיכך נוכל גם לקבוע את

\medskip
\noindent\underline{משפט 14.} כאשר המספר $ab$ הנו אחת מהצורות
\[ K^2 + 1 \]
\[ K'^2 + 4 \qquad \text{או} \]
בשביל מספר $K'$ אי-זוגי, בתנאי
\[ K' > 1 \]
אזי אין פתרון למשוואה (1).

\noindent ואמנם, כידוע, קיים
\[ \sqrt{K^2+1} = (K, \overline{2K}) \]
וכן
\[ \sqrt{K^2+4} = \left(K, \overline{\tfrac{K-1}{2}, 1, 1, \tfrac{K-1}{2}, 2K}\right) \]
בשביל כל $K$ אי-זוגי $>1$.

\noindent מאידך נוכל, כמסקנה מן התוצאות שקבלנו, לקבוע את

\medskip
\noindent\underline{משפט 15.} המחזור של השבר המשולב $\sqrt{K(K\pm1)}$, בשביל $K>2$, מרכב ממספר זוגי של אברים.

\noindent לשם הוכחה מספיק לשים לב לשיוויון

% [original page: 82]
\[ (K \pm 1)x^2 - Ky^2 = \pm 1 \]
המתקיים בשביל
\[ x = y = 1 . \]
אולם לפי זה נוכל גם להסיק את

\medskip
\noindent\underline{משפט 16.} למשוואה
\[ ax^2 - by^2 = \pm 1 \]
אין פתרון בשביל
\[ ab = K(K\pm1) \]
כאשר
\[ \left. \begin{aligned} 1 &< a \neq K\pm1 \\ 1 &< b \neq K \end{aligned} \right\} \]
מסקנה זאת מתקבלת מיד על סמך משפט 11 בשימת לב לשיוויון
\[ (K\pm1)x^2 - Ky^2 = \pm 1 \]
בשביל
\[ x = y = 1 \]

\bigskip
\noindent\textbf{4.} נראה עתה שלגבי פתרונות המשוואה (1) קיימת נוסחה דומה לזאת הקיימת לגבי משוואת פל בשינוי-מה.

\noindent לפי נוסחאות (12) נוכל לכתוב, בהסתמך על הנוסחאות הידועות הקיימות לגבי משוואת פל
\[ \mathcal{E}^2 - ab\eta^2 = 1, \qquad ax_n^2 + by_n^2 + 2x_ny_n\sqrt{ab} = (\mathcal{E}_1 + \eta_1\sqrt{ab})^{2n-1} \]
בשביל כל
\[ n = 1 , 2 , 3 , \ldots \]
ומכאן: אבל
\[ (x_n\sqrt{a} + y_n\sqrt{b})^2 = (\mathcal{E}_1 + \eta_1\sqrt{ab})^{2n-1} \]
\[ (\mathcal{E}_1 + \eta_1\sqrt{ab})^1 = (x_1\sqrt{a} + y_1\sqrt{b})^2 \]
ולכן
\[ (x_n\sqrt{a} + y_n\sqrt{b})^2 = (x_1\sqrt{a} + y_1\sqrt{b})^{2(2n-1)} \]
ומכאן נובע כי:
\[ x_n\sqrt{a} + y_n\sqrt{b} = (x_1\sqrt{a} + y_1\sqrt{b})^{2n-1} \]
בשביל
\[ n = 1 , 2 , 3 , \ldots \]
נביא עתה בחשבון גם את המקרה
\[ \eta_1 \equiv 1 \pmod{2} \]

% [original page: 83]
\noindent שבו קיים משפט 10, ואז יתקיים
\[ (x_n\sqrt{a} + y_n\sqrt{b})^2 = (\mathcal{E}_1 + \eta_1\sqrt{ab})^{2(2n-1)} = (\mathcal{E}_2 + \eta_2\sqrt{ab})^{2n-1} \]
ולפיכך נקבל שוב בהתאם למשפט 10:
\[ x_n\sqrt{a} + y_n\sqrt{b} = (x_1\sqrt{a} + y_1\sqrt{b})^{2n-1} . \]
נרחיב עתה את הנוסחה לכל קומבינציות הסימנים האפשריות
\[ (\pm x_n , \pm y_n) \]
ונקבל לבסוף את הנוסחה:
\[ \boxed{x_n\sqrt{a} + y_n\sqrt{b} = \pm(x_1\sqrt{a} \pm y_1\sqrt{b})^{2n-1}} \tag{13} \]
בשביל כל $n = 1 , 2 , 3 , \ldots$
% NOTE: the right-hand subscripts in the boxed formula (13) are printed ambiguously in the source (possibly "x_n" rather than "x_1"); rendered here as x_1, y_1 for internal consistency with the immediately preceding derivation. See transcription notes.

\medskip
\noindent\textbf{\underline{הערה.}} המשוואה
\[ u^2 - Dv^2 = -1 \]
מהווה מקרה מיוחד של המשוואה (1) הכללית, כאשר
\[ e = -1 \]
או בצורה
\[ Dv^2 - u^2 = 1 . \]

% [original page: 84]
\subsection*{\S3. הכללה שניה.}

\noindent\textbf{1.} כהכללה שניה של משוואת פל אנו רואים במשוואה מהצורה
\[ ax^2 - by^2 = \pm 2 \tag{1} \]
כאשר $a,b$ הנם מספרים אי-זוגיים. כפי שקל לראות יהיו שני המספרים
\[ ax^2 - by^2 \, , \quad ax^2 + by^2 \]
זוגיים, ולכן מותר להניח:
\[ \begin{cases} ax_n^2 + by_n^2 = 2\mathcal{E} \\ x_ny_n = \eta \end{cases} \tag{2} \]
בשביל $\mathcal{E}$, $\eta$ שלמים מסויימים.

\noindent אולם, כפי שאפשר להיווכח מיד מתקיים בהחלט זאת
\[ \mathcal{E}^2 - ab\eta^2 = 1 \, ; \tag{3} \]
כלומר המספרים $\mathcal{E}$, $\eta$ מקיימים את משוואת פל הידועה.

\noindent בדרך דומה לזאת שהלכנו בה בסעיף הקודם ניווכח, שהתנאי ההכרחי והמספיק לקיום הקשרים (1)(2) הוא שתתקיימנה הקונגרואנציות
\[ \mathcal{E} \pm 1 \equiv 0 \pmod{a} \tag{4} \]
\[ \mathcal{E} \mp 1 \equiv 0 \pmod{b} \]
ונוסף עליהן:
\[ \mathcal{E} \equiv 0 \pmod{2} \]
אולם מכיוון ש:
\[ xy \equiv ab \equiv 1 \pmod{2} \]
נובע מיד, כי $\eta$ הנו אי-זוגי, לפי (2), וכן נסיק מ-(4) שמתקיימים
\[ \mathcal{E} + 1 \equiv \mathcal{E} - 1 \equiv 1 \pmod{2} . \]
בדרך דומה לזאת שבסעיף הקודם ניווכח שהקונגרואנציות
\[ \mathcal{E}_{2n-1} \pm 1 \equiv 0 \pmod{a} \]
\[ \mathcal{E}_{2n-1} \mp 1 \equiv 0 \pmod{b} \]
קיימות בשביל כל $n$ אם הן קיימות בשביל $n$ מסויים, בתנאי הנוסף
\[ \mathcal{E}_n + 1 \equiv \mathcal{E}_n - 1 \equiv 1 \pmod{2} \]

% [original page: 85]
\noindent לפי"ז אנו מקבלים, איפוא, את

\medskip
\noindent\underline{משפט 17.} כל פתרונות המשוואה (1) נכללים בנוסחאות
\[ \begin{cases} ax_n^2 + by_n^2 = 2\mathcal{E}_{2n-1} \\ x_ny_n = \eta_{2n-1} \end{cases} \qquad n = 1 , 2 , 3 , \ldots \tag{5} \]
אם יש בכלל פתרון למשוואה הנ"ל.

\noindent מתוך המסקנות שבסעיפים הקודמים אנו מקבלים מיד את

\medskip
\noindent\underline{משפט 18.} אם יש פתרון למשוואה
\[ ax^2 - by^2 = \pm 1 \]
אזי אין פתרון למשוואה
\[ ax^2 - by^2 = \pm 2 . \]
ולהיפך.

\noindent עובדה זאת נובעת מהדרישה ש-$\eta_1$ יהיה זוגי בשביל קיום המשוואה
\[ ax^2 - by^2 = \pm 1 \]
ובשביל קיום המשוואה
\[ ax^2 - by^2 = \pm 2 \]
דרוש ש-$\eta_1$ יהיה לא זוגי.

\medskip
\noindent\textbf{\underline{הערה.}} המשפטים 11, 12, 13, 14 נשארים בתקפם גם לגבי המשוואה:
\[ ax^2 - by^2 = \pm 2 \]

\medskip
\noindent\underline{משפט 19.} אם אפשר לפרק פונקציה $F(K)$ ממספר שלם לשני גורמים מספרים שלמים $f(K), g(K)$ באופן שמתקיים אחד השווויונים
\[ f(K)x^2 - g(K)y^2 = \pm 1 , \pm 2 \]
בשביל זוג מספרים שלמים $(x,y)$ אז מרכב המחזור שבפיתוח $\sqrt{F(K)}$ לשבר משולב ממספר זוגי של אברים.

\noindent כי במקרה ההפוך היה מתקיים

% [original page: 86]
\[ F(K)x^2 - y^2 = 1 \]
ולכן לא היה פתרון למשוואות האחרות.

\noindent לדוגמה, כאשר
\[ F(K) = K(K\pm2); \]
\[ K(K\pm1); \quad K>2; \]
\[ n(K^2n\pm2); \]
\[ 4K^2 \pm 25K + 39 \quad (n>4) \]
אזי מרכב מחזור השבר המשולב $\sqrt{F(K)}$ ממספר אברים זוגי, מכיוון שהסווררונים:
\[ (K\pm2) \cdot 1^2 - K \cdot 1^2 = \pm 2 , \]
\[ (K^2n\pm2) \cdot 1^2 - nk^2 = \pm 2 , \]
\[ (K\pm1) \cdot 1^2 - K \cdot 1^2 = \pm 1 \]
\[ (4n\pm13) \cdot 1^2 - (n\pm3) \cdot 2^2 = \pm 1 , \]
מתקיימים באופן זהותי.

\bigskip
\noindent\textbf{2.} נראה עתה שגם לגבי המשוואה
\[ ax^2 - by^2 = \pm 2 \]
קיימת נוסחה למציאת כל הפתרונות בעזרת הפתרון המינימלי הדומה לקודמות.

\noindent לפי הנוסחאות (5) נקבל על סמך המשפטים היסודיים הקיימים לגבי משוואת פל:
\[ \left(\frac{x_n\sqrt{a} + y_n\sqrt{b}}{\sqrt{2}}\right)^2 = (\mathcal{E}_1 + \eta_1\sqrt{ab})^{2n-1} \]
ומכאן:
\[ \left(\frac{x_n\sqrt{a} + y_n\sqrt{b}}{\sqrt{2}}\right)^2 = \left(\frac{x_1\sqrt{a} + y_1\sqrt{b}}{\sqrt{2}}\right)^{2(2n-1)} \]
ולכן נקבל:
\[ \boxed{\frac{x_n\sqrt{a} + y_n\sqrt{b}}{\sqrt{2}} = \pm\left(\frac{x_1\sqrt{a} \pm y_1\sqrt{b}}{\sqrt{2}}\right)^{2n-1}} \tag{6} \]
בשביל כל $n = 1, 2, 3, \ldots$

% [original page: 87]
\noindent\textbf{3.} לבסוף נוכיח עוד את

\medskip
\noindent\underline{משפט 20.} בשביל כל מספר $D$ שלם יש פתרון בשביל אחת המשוואות
\[ ax^2 - by^2 = 1 \]
או
\[ ax^2 - by^2 = -1 \]
בתנאי
\[ a>1 \]
כאשר
\[ ab = D. \]
הנו פירוק מסויים של $D$ לגורמים זרים ביניהם.

\medskip
\noindent\textbf{\underline{הוכחה.}}
נוכל להיווכח מהשוויון
\[ (\mathcal{E}_1 + 1)(\mathcal{E}_1 - 1) = ab\eta_1^2 \]
הקיים תמיד; לו היה
\[ \left. \begin{aligned} \mathcal{E}_1 - 1 &= 2D\eta'^2 \\ \mathcal{E}_1 + 1 &= 2\mathcal{E}'^2 \end{aligned} \right\} \]
היה נובע מיד, כי:
\[ \mathcal{E}'^2 - D\eta'^2 = 1 \]
בשביל
\[ \eta'\mathcal{E}' = \eta_1 \]
כאשר
\[ \mathcal{E}' > 1 \]
מוכרח להיות
\[ \eta' < \eta_1 . \]
וזה מתנגד להנחה ש-$(\mathcal{E}_1, \eta_1)$ הנו הפתרון המינימלי של משוואת פל.

% [original page: 88]
\subsection*{\S4. הכללה שלישית.}

\noindent\textbf{1.} אנו עוברים עתה להכללה שלישית של משוואת פל, והיא למשוואה
\[ ax^2 - by^2 = \pm 4 \tag{1} \]
בתנאים:
\[ \left. \begin{aligned} (a,b) &= 1 \\ (x,y) &= 1 \end{aligned} \right\} \tag{2} \]
בראש וראשונה נוכיח את

\medskip
\noindent\underline{משפט 21.} אם יש למשוואה (1) פתרונות בתנאי $(x,y)=1$ אזי יש לה גם פתרונות לא-עיקריים: במלים אחרות: מקיום המשוואה (1) נובע קיום המשוואה
\[ a\alpha^2 - b\gamma^2 = \pm 1 \]

\medskip
\noindent\textbf{\underline{הוכחה.}}
נסתמך כאן על הזהות:
\[ (x\sqrt{a} + y\sqrt{b})^3 = x\sqrt{a}(ax^2+3by^2) + y\sqrt{b}(3ax^2+by^2) \]
ונקבל לפי (1):
\[ (x\sqrt{a} + y\sqrt{b})^3 = 8x\sqrt{a}\left(\frac{by^2\pm1}{2}\right) + 8y\sqrt{b}\left(\frac{ax^2\mp1}{2}\right) \]
המספרים $ax^2\mp1, by^2\pm1$ הם שניהם בבת אחת זוגיים, כפי שנובע מ-(1), ולפיכך נקבל:
\[ (ax^2 - by^2)^3 = 8^2(as^2 - bt^2) \tag{3} \]
בשביל מספרים שלמים $s,t$ מסויימים. אולם מכאן נובע מיד, לפי (1), כי:
\[ 8^2(as^2 - bt^2) = (\pm4)^3 \]
ולכן
\[ as^2 - bt^2 = \pm 1 \tag{4} \]
ומתקיימים הקשרים
\[ \left. \begin{aligned} s &= x\left(\frac{by^2\pm1}{2}\right) \\ t &= y\left(\frac{ax^2\mp1}{2}\right) \end{aligned} \right\} \tag{5} \]
\hfill מ.ש.ל.

\medskip
\noindent\textbf{2.} לפי ההנחה $ab$ הנו אי-זוגי אפשר תמיד להניח:

% [original page: 89]
\[ \begin{cases} ax^2 + by^2 = 2u \\ xy = v \end{cases} \tag{6} \]
בשביל מספרים $u,v$ שלמים, מכיוון שגם המספרים $x,y$ הם בבת-אחת זוגיים או אי-זוגיים.

\noindent אבל לפי"ז נקבל ע"י העלאת שני אגפי המשוואה (1) לריבוע:
\[ u^2 - Dv^2 = 4 , \quad (D=ab) \tag{7} \]
והמספרים $u,v$ הם בבת-אחת זוגיים או אי-זוגיים.

\noindent לפיכך מתקבל

\medskip
\noindent\underline{משפט 22.} כאשר $ab$ הנו מהצורה $8K+1$ אזי אין פתרון למשוואה (1).

\noindent ואמנם, אין פתרון למשוואה
\[ u^2 - (8K+1)v^2 = 4 \]
בתנאי
\[ (u,v) = 1 \]
מכיוון, שאז
\[ u^2 \equiv v^2 \equiv 1 \pmod{8} \]
ולכן
\[ u^2 - v^2 \equiv 0 \pmod{8} \]
בו בזמן ש:-
\[ 8K^2 + 4 = 4(2Kv^2+1) \not\equiv 0 \pmod{8} \]

\medskip
\noindent\textbf{3.} נוכיח עתה משפט הפוך בתנאי מסויים למשפט 21:

\medskip
\noindent\underline{משפט 23.} אם יש פתרון למשוואה
\[ a\alpha^2 - b\gamma^2 = \pm 1 \]
וגם למשוואת פל
\[ u^2 - abv^2 = 4 \]
בתנאי
\[ (u,v) = 1 \]
אזי יש פתרון למשוואה
\[ ax^2 - by^2 = \pm 4 \]
בתנאי
\[ (x,y) = 1 \]

\medskip
\noindent\textbf{\underline{הוכחה.}}
שתי התבניות הריבועיות $(a,0,-b)$, $(1,0,-ab)$ הן

% [original page: 90]
\noindent אקוויולנטות בעזרת הצגת
\[ \begin{pmatrix} \alpha_1 & b\beta_1 \\ \beta_1 & a\alpha_1 \end{pmatrix} \]
בעלת הדטרמיננטה
\[ a\alpha_1^2 - b\beta_1^2 = \pm 1 \]
ומכאן נכונות המשפט.

\medskip
\noindent\textbf{4.} באופן דומה לזה שבסעיפים הקודמים ניווכח שהתנאי ההכרחי בשביל קיום המערכות (6)(7) כאשר $u$ הנו אי-זוגי הוא קיום הקונגרואנציות
\[ u \pm 2 \equiv 0 \pmod{a} \]
\[ u \mp 2 \equiv 0 \pmod{b} \tag{8} \]
אולם על סמך משפט 21 ניווכח, שהקונגרואנציות הנ"ל הן גם מספיקות בשביל $u$ זוגי, אם הן מתקיימות בשביל $u$ אי-זוגי.

\noindent כי בתנאים אלו קיימת המשוואה
\[ as^2 - bt^2 = \pm 1 \]
ולכן, לפי המשפטים שבסעיף 2, מוכרח להיות $\mathcal{E}$ המקיים את המשוואה
\[ \mathcal{E}^2 - ab\eta^2 = 1 \]
אי-זוגי.

\noindent מכאן, שהמספרים,
\[ u \pm 2 = 2(\mathcal{E}\pm1) \]
\[ u \mp 2 = 2(\mathcal{E}\mp1) \]
מתחלקים שניהם ב-4, נשים עוד לב לכך, כי:
\[ (u+2 \, , \, u-2) \mid 4 \]
ונסיק מן השוויון:
\[ (u-2)(u+2) = ab \cdot v^2 \]
בהתאם לקונגרואנציות (8), כי:
\[ u \pm 2 = 4ap^2 \]
\[ u \mp 2 = 4bq^2 \]
בשביל מספרים $p,q$ שלמים, וזה מתאים למערכת הקשרים (6).

% [original page: 91]
\noindent\textbf{5.} נראה עתה שאפשר לקבל נוסחאות כלליות למשוואה הנכללת על דוגמת הנוסחאות הקיימות למשוואת פל עצמה.

\noindent לגבי משוואת פל עצמה קיימות הנוסחאות הכלליות:
\[ 2^{n-1}(u_n + v_n\sqrt{ab}) = (u_1 + v_1\sqrt{ab})^n \tag{9} \]
בתנאים הידועים.

\noindent מכאן נובע מיד, כי:
\[ 2^{n-1}u_n = u_1^n + ab \cdot t \]
כש-$t$ הנו מספר שלם וחיובי מסויים.

\noindent כאשר $n$ הוא מספר אי-זוגי מתקיימים, איפוא, הקשרים
\[ 2^{n-1}(u_n+2) = (u_1+2)P + ab \cdot t \]
\[ 2^{n-1}(u_n-2) = (u_1-2)Q + ab \cdot t \]
בשביל $P,Q$ שלמים מסויימים.

\noindent מכאן רואים מיד, שאם מתקיימות הקונגרואנציות (8) בשביל $u_1$ הרי הן מתקיימות גם בשביל כל $u_n$ עם $n$ אי-זוגי, כאשר $a,b$ הם מספרים אי-זוגיים. אולם בדרך דומה לזאת שהוכחנו ביחס להכללות הקודמות ניווכח בנקל, שגם ההיפוך נכון, כלומר: אם הקונגרואנציות הנ"ל מתקיימות לגבי $u_n$ הרי הן מתקיימות גם לגבי $u_1$. חוץ מזה אין הן יכולות להתקיים לגבי $u_n$ עם $n$ זוגי. מכאן המסקנה:

\noindent כל הפתרונות $x_n, y_n$ של המשוואה (1) נכללים בנוסחאות:
\[ ax_n^2 + by_n^2 = 2u_{2n-1} \, , \quad x_ny_n = v_{2n-1} \tag{10} \]
בשביל כל $n$ טבעי; $n \geq 1$.

\noindent על סמך הנוסחה הסימבולית (9) נקבל איפוא מיד:
\[ \left(\frac{x_n\sqrt{a} + y_n\sqrt{b}}{2}\right)^2 = \left(\frac{x_1\sqrt{a} + y_1\sqrt{b}}{2}\right)^{(2n-1)2} \]

% === TRANSCRIPTION NOTES ===
% - Page 75 (p083): one sentence about restricting to natural values of x,y,E,eta by symmetry is partly obscured by scan noise/dot artifacts; transcribed to the best approximation of visible words. Flagged with an inline ILLEGIBLE comment at that location.
% - Page 76 (p084): the source shows a small hand-drawn arrow mark (pen annotation, "^") near the displayed congruences; treated as marginalia per the style guide and omitted from the transcription.
% - Page 77 (p085) and page 78 (p086): the source contains fraction-like stacked notations ("u" over "\eta_{2n}", "p_n" over "2\eta_{2n-1}", "q_n" over "2\eta_{2n-1}") that are ambiguous between a genuine typeset fraction and a divisibility annotation. Interpreted as divisibility statements (u \mid \eta_{2n}, p_n \mid 2\eta_{2n-1}, q_n \mid 2\eta_{2n-1}) per the style guide's convention for such number-theoretic stacked layouts; the exact intended meaning could not be fully confirmed from the scan quality.
% - Page 78 (p086): the Dickson citation ("History of the theory of numbers", vol. II, p. 392) and the quoted theorem statement directly below it are transcribed from a somewhat smudged section of the scan; wording of the quoted paraphrase is a best-effort reconstruction of the visible Hebrew text.
% - Page 83 (p091): boxed formula (13) — the right-hand side subscripts of x and y are unclear in the scan (could read as x_n/y_n or x_1/y_1); rendered as x_1, y_1 for consistency with the surrounding derivation, which explicitly builds from the minimal solution (x_1, y_1). Flagged inline.
% - Throughout: the cursive/script letters resembling stylized "E" and a cursive lowercase letter were rendered as \mathcal{E} and \eta respectively, consistent with the golden example and with the plain "eta" glyph appearing directly in some printed formulas (e.g. p.75, p.77).
% - Bold/underlined structural labels (הוכחה, הערה, משפט N) were rendered per the style guide's \textbf{\underline{...}} convention; the bare paragraph numbers ("1.", "2.", "3." ...) that mark sub-parts within a section (not full TOC-style headers) were rendered as \textbf{N.} inline rather than as \subsection* headers, since they introduce a paragraph rather than a titled section - only the two section headers "§3. הכללה שניה" and "§4. הכללה שלישית" were rendered as \subsection*.
% - Page 86 (p094): a few short inline expressions in the worked examples of Theorem 19 (e.g. "n(K^2n \pm 2)", "(K^2n\pm2)\cdot1^2 - nk^2") are transcribed exactly as printed, including what appear to be inconsistent capitalization of the variable (K vs k) in the original source; this is preserved for fidelity rather than corrected.
